\documentclass[11pt]{article}
\usepackage[T1]{fontenc}
\usepackage{lmodern}
\usepackage[margin=1in]{geometry}
\usepackage{amsmath,amssymb,amsthm,mathtools}
\usepackage[hidelinks]{hyperref}
\usepackage{enumitem}
\emergencystretch=1.5em

\newtheorem{theorem}{Theorem}
\newtheorem{lemma}[theorem]{Lemma}
\newtheorem{corollary}[theorem]{Corollary}
\newtheorem{remark}[theorem]{Remark}
\newcommand{\R}{\mathbb R}
\newcommand{\B}{B}
\newcommand{\J}{J}
\newcommand{\DOH}{\mathrm{DOH}}
\newcommand{\wSSD}{w^*\text{-}\mathrm{SSD2P}}

\title{Decomposable Octahedrality Is Dual to the\\
Weak-Star Symmetric Strong Diameter-Two Property}
\author{Run \texttt{fa\_banach\_001}, lane 14}
\date{August 9, 2026}

\begin{document}
\maketitle

\begin{abstract}
Ostrak proved that if $X^*$ has the weak-star symmetric strong
diameter-two property, then $X$ is decomposably octahedral, and asked
whether the converse holds for every Banach space.  We prove the converse.
The simultaneous dual-ball constraints are encoded as the unit ball of the
norm
\[
 (f_1,\ldots,f_n,g)\longmapsto
 \max_i\bigl\{\|f_i+g\|,\|f_i-g\|\bigr\}.
\]
Hahn--Banach expresses its support function by decompositions in $X^{**}$.
The principle of local reflexivity transfers each finite decomposition to
$X$, where decomposable octahedrality supplies exactly the required lower
bound.  A finite-dimensional separation argument then produces one large
functional $g$ that works simultaneously in all prescribed weak-star
slices.  Consequently, for every real Banach space $X$, $X$ is
decomposably octahedral if and only if $X^*$ has the weak-star symmetric
strong diameter-two property.
\end{abstract}

\section{The question and the result}

All Banach spaces are real.  A weak-star slice of $B_{X^*}$ is
\[
 S(B_{X^*},x,\alpha)
 =\{f\in B_{X^*}:f(x)>1-\alpha\},
 \qquad x\in S_X,\quad \alpha>0.
\]
The dual $X^*$ has the \emph{weak-star symmetric strong diameter-two
property} ($\wSSD$) if, for every finite family of weak-star slices
$S_1,\ldots,S_n$ and every $\varepsilon>0$, there are $f_i\in S_i$ and
$g\in B_{X^*}$ such that
\[
 f_i+g\in S_i,\qquad f_i-g\in S_i\quad(1\le i\le n),
 \qquad \|g\|>1-\varepsilon.
\]

Following Ostrak \cite{Ostrak}, $X$ is \emph{decomposably octahedral}
($\DOH$) if, for every finite $E\subset S_X$ and every $\rho>0$, there is
$y\in S_X$ such that, whenever
\[
 y=\sum_{i=1}^m y_i,\qquad x_i\in E,\qquad a_i,b_i\ge0,
\]
one has
\begin{equation}\label{eq:doh}
 \sum_{i=1}^m
 \bigl(\|a_i x_i+y_i\|+\|b_i x_i-y_i\|\bigr)
 \ge (1-\rho)
 \left(\sum_{i=1}^m(a_i+b_i)+2\right).
\end{equation}
Ostrak proved that $\wSSD$ of $X^*$ implies $\DOH$ of $X$ and asked
whether the converse holds \cite[Proposition~2.1 and p.~7]{Ostrak}.
The question was still recorded as open in 2024
\cite[Section~4, paragraph preceding Question~4.2]{HKO}.

\begin{theorem}[Full converse]\label{thm:main}
Let $X$ be a real Banach space.  If $X$ is decomposably octahedral, then
$X^*$ has the weak-star symmetric strong diameter-two property.
\end{theorem}

Combined with Ostrak's proposition, this gives the exact duality.

\begin{corollary}\label{cor:equivalence}
For every real Banach space $X$, the following are equivalent:
\begin{enumerate}[label=\textup{(\roman*)}]
\item $X$ is decomposably octahedral;
\item $X^*$ has the weak-star symmetric strong diameter-two property.
\end{enumerate}
\end{corollary}

\section{Proof intuition}

Fix slice directions $x_1,\ldots,x_n$.  We seek functionals $f_i$ and a
common $g$ for which both $f_i+g$ and $f_i-g$ lie in the dual unit ball and
are almost norming at $x_i$.  Add one more desired coordinate, $g(y)$,
where $y$ is supplied by $\DOH$.  If the resulting compact convex set of
$2n+1$ coordinates missed the upper orthant near $(1,\ldots,1)$, a
nonnegative linear functional would separate them.

It therefore suffices to lower-bound every nonnegative support functional.
The constraints $\|f_i\pm g\|\le1$ define the unit ball of an equivalent
norm on $(X^*)^{n+1}$.  Hahn--Banach computes its dual norm as a minimum of
\[
 \sum_i\bigl(\|a_i\J x_i+r_i\|+
                    \|b_i\J x_i-r_i\|\bigr),
 \qquad 2\sum_i r_i=c\J y,
\]
over $r_i\in X^{**}$.  Formula \eqref{eq:doh} is precisely the needed
lower bound when the $r_i$ belong to $X$.  Local reflexivity sends any
finite family in $X^{**}$ almost isometrically into $X$, fixes the displayed
vectors from $X$, and preserves the sum.  Letting the local-reflexivity
distortion tend to one proves the same bound in $X^{**}$.  This closes the
duality gap and makes separation impossible.

\section{The support-function formula}

Write $\J:X\to X^{**}$ for the canonical embedding.

\begin{lemma}[Hahn--Banach support formula]\label{lem:support}
Fix $x_1,\ldots,x_n,y\in X$ and define
\[
 \mathcal C=
 \left\{(f_1,\ldots,f_n,g)\in(X^*)^{n+1}:
       \|f_i+g\|\le1,\ \|f_i-g\|\le1\ (1\le i\le n)\right\}.
\]
For $a_i,b_i,c\ge0$, let
\begin{align*}
 H(a,b,c)=\sup_{(f,g)\in\mathcal C}
 \bigg[&\sum_{i=1}^n a_i(f_i+g)(x_i)
       +\sum_{i=1}^n b_i(f_i-g)(x_i)+c g(y)\bigg].
\end{align*}
Then
\begin{equation}\label{eq:support}
 H(a,b,c)=
 \min_{\substack{r_1,\ldots,r_n\in X^{**}\\
                   2\sum_i r_i=c\J y}}
 \sum_{i=1}^n
 \bigl(\|a_i\J x_i+r_i\|+\|b_i\J x_i-r_i\|\bigr).
\end{equation}
\end{lemma}

\begin{proof}
On $Z=(X^*)^{n+1}$ define
\[
 N(f_1,\ldots,f_n,g)
 =\max_{1\le i\le n}\{\|f_i+g\|,\|f_i-g\|\}.
\]
The map
\[
 A:Z\longrightarrow W=(X^*)^{2n},\qquad
 A(f_1,\ldots,f_n,g)=(f_i+g,f_i-g)_{i=1}^n,
\]
is an isometry when $W$ has its maximum norm; its range is closed because
$f_i$ and $g$ are recovered by taking half-sums and half-differences.
Thus $\mathcal C$ is the $N$-unit ball.

The displayed functional defining $H(a,b,c)$ is an element of $Z^*$.
Since the $N$-unit ball is symmetric, $H(a,b,c)$ is its dual norm.
By Hahn--Banach, this norm is the least norm of an extension from
$A(Z)$ to $W$.  Since $W^*$ is the $\ell_1$-sum of $2n$ copies of $X^{**}$,
such an extension is represented by $p_i,q_i\in X^{**}$ and has norm
$\sum_i(\|p_i\|+\|q_i\|)$.  The restriction agrees with the displayed
functional exactly when
\begin{align}
 p_i+q_i&=(a_i+b_i)\J x_i,\label{eq:adjoint1}\\
 \sum_{i=1}^n(p_i-q_i)
 &=\sum_{i=1}^n(a_i-b_i)\J x_i+c\J y.\label{eq:adjoint2}
\end{align}
Put $r_i=p_i-a_i\J x_i$.  Equations
\eqref{eq:adjoint1}--\eqref{eq:adjoint2} become
\[
 p_i=a_i\J x_i+r_i,\qquad q_i=b_i\J x_i-r_i,
 \qquad 2\sum_i r_i=c\J y.
\]
Conversely, every such family $(r_i)$ gives an extension.  Hahn--Banach
provides an extension of exactly the dual norm, so the minimum is attained
and \eqref{eq:support} follows.
\end{proof}

\section{DOH controls the bidual decompositions}

\begin{lemma}[Local-reflexivity transfer]\label{lem:lr}
Assume that $X$ is $\DOH$.  Fix a finite set
$E=\{x_1,\ldots,x_n\}\subset S_X$ and $\rho>0$, and choose $y\in S_X$
as in \eqref{eq:doh}.  Then, for every $a_i,b_i,c\ge0$,
\begin{equation}\label{eq:lower-support}
 H(a,b,c)\ge
 (1-\rho)\left(\sum_{i=1}^n(a_i+b_i)+c\right),
\end{equation}
where $H$ is defined in Lemma~\ref{lem:support}.
\end{lemma}

\begin{proof}
By Lemma~\ref{lem:support}, it is enough to bound the objective in
\eqref{eq:support} for an arbitrary admissible family $(r_i)$.

If $c=0$, then $\sum_i r_i=0$, although even this equality is unnecessary:
the triangle inequality gives, term by term,
\[
 \|a_i\J x_i+r_i\|+\|b_i\J x_i-r_i\|
 \ge(a_i+b_i)\|\J x_i\|=a_i+b_i.
\]

Suppose that $c>0$.  Let $F\subset X^{**}$ be the finite-dimensional span
of $r_1,\ldots,r_n,\J x_1,\ldots,\J x_n,\J y$.  The principle of local
reflexivity \cite{LR} says that, for every $\tau>0$, there is a linear map
$T_\tau:F\to X$ such that
\[
 T_\tau(\J x)=x\quad(\J x\in F\cap\J X),
 \qquad \|T_\tau\|\le1+\tau.
\]
Set $y_i=(2/c)T_\tau r_i$.  Since $2\sum_i r_i=c\J y$ and
$T_\tau(\J y)=y$, we have $\sum_i y_i=y$.  Applying \eqref{eq:doh} gives
\begin{align*}
 &\sum_i\bigl(\|a_i x_i+T_\tau r_i\|
                  +\|b_i x_i-T_\tau r_i\|\bigr)\\
 &\qquad=\frac c2\sum_i
 \left(\left\|\frac{2a_i}{c}x_i+y_i\right\|
       +\left\|\frac{2b_i}{c}x_i-y_i\right\|\right)\\
 &\qquad\ge(1-\rho)\left(\sum_i(a_i+b_i)+c\right).
\end{align*}
On the other hand, because $T_\tau$ fixes each $\J x_i$,
\begin{align*}
 &\sum_i\bigl(\|a_i\J x_i+r_i\|+
                   \|b_i\J x_i-r_i\|\bigr)\\
 &\qquad\ge\frac1{1+\tau}
 \sum_i\bigl(\|a_i x_i+T_\tau r_i\|
                  +\|b_i x_i-T_\tau r_i\|\bigr).
\end{align*}
Letting $\tau\downarrow0$ proves the desired lower bound for every
admissible $(r_i)$.  Formula \eqref{eq:support} now yields
\eqref{eq:lower-support}.
\end{proof}

\section{Separation and completion of the proof}

\begin{proof}[Proof of Theorem~\ref{thm:main}]
Let
\[
 S_i=S(B_{X^*},x_i,\alpha_i),\qquad 1\le i\le n,
\]
be weak-star slices, and fix $\varepsilon>0$.  Choose
\[
 0<\rho<\min\{\varepsilon,\alpha_1,\ldots,\alpha_n\},
\]
and let $y\in S_X$ be supplied by $\DOH$ for
$E=\{x_1,\ldots,x_n\}$ and the parameter $\rho$.

For the set $\mathcal C$ in Lemma~\ref{lem:support}, consider the compact
convex subset $K\subset\R^{2n+1}$ consisting of the vectors
\[
 \bigl((f_i+g)(x_i),(f_i-g)(x_i)\bigr)_{i=1}^n,\ g(y),
 \qquad (f_1,\ldots,f_n,g)\in\mathcal C.
\]
Indeed, $\mathcal C$ is weak-star closed and lies in the product of
$n+1$ dual unit balls: both $f_i$ and $g$ are half-sums or
half-differences of $f_i+g$ and $f_i-g$.  Banach--Alaoglu therefore gives
the asserted compactness.

We claim that
\begin{equation}\label{eq:orthant}
 K\cap[1-\rho,\infty)^{2n+1}\ne\varnothing.
\end{equation}
Otherwise, finite-dimensional strong separation of the compact set $K$
from the closed upper orthant gives coefficients
$a_i,b_i,c\ge0$, not all zero, such that
\begin{align*}
 &\sup_{(f,g)\in\mathcal C}
 \left[\sum_i a_i(f_i+g)(x_i)
      +\sum_i b_i(f_i-g)(x_i)+c g(y)\right]\\
 &\hspace{5cm}<
 (1-\rho)\left(\sum_i(a_i+b_i)+c\right).
\end{align*}
The coefficients are nonnegative because a linear functional bounded
below on the upper orthant must have nonnegative coordinates.  The strict
inequality contradicts Lemma~\ref{lem:lr}.  This proves
\eqref{eq:orthant}.

Choose $(f_1,\ldots,f_n,g)\in\mathcal C$ producing a point in the
intersection.  Then
\[
 \|f_i\pm g\|\le1,\qquad
 (f_i\pm g)(x_i)\ge1-\rho>1-\alpha_i.
\]
Thus $f_i+g,f_i-g\in S_i$.  Their midpoint $f_i$ also belongs to the
convex slice $S_i$.  Finally,
\[
 \|g\|\ge g(y)\ge1-\rho>1-\varepsilon.
\]
These are exactly the requirements of the $\wSSD$.
\end{proof}

\begin{proof}[Proof of Corollary~\ref{cor:equivalence}]
The implication \textup{(ii)}$\Rightarrow$\textup{(i)} is
\cite[Proposition~2.1]{Ostrak}; the reverse implication is
Theorem~\ref{thm:main}.
\end{proof}

\section{Audit note}

The point at which a naive duality proof can appear to fail is important.
Hahn--Banach naturally produces $r_i\in X^{**}$, whereas the definition of
$\DOH$ quantifies only over decompositions in $X$.  The finite-dimensional
span of a single candidate decomposition is precisely the setting of local
reflexivity.  Because the sum of the $r_i$ is the canonical image of the
fixed vector $cy/2$, a local-reflexivity map both preserves that sum and
fixes every $x_i$.  No approximation of infinitely many vectors and no
compactness passage is needed.

\begin{thebibliography}{9}

\bibitem{Ostrak}
A. Ostrak,
\emph{On the duality of the symmetric strong diameter $2$ property in
Lipschitz spaces},
Rev. R. Acad. Cienc. Exactas F\'is. Nat. Ser. A Mat. RACSAM
\textbf{115} (2021), Paper No.~66;
arXiv:2008.03163.

\bibitem{HKO}
R. Haller, J. K. Kaasik, and A. Ostrak,
\emph{Separating diameter two properties from their weak-star counterparts
in spaces of Lipschitz functions},
arXiv:2404.11430.

\bibitem{LR}
J. Lindenstrauss and H. P. Rosenthal,
\emph{The $\mathcal L_p$ spaces},
Israel J. Math. \textbf{7} (1969), 325--349.

\end{thebibliography}

\end{document}
